\documentclass[12pt]{article}
\usepackage[utf8]{inputenc}
\usepackage[T1]{fontenc}
\usepackage[margin=1in]{geometry}
\usepackage{setspace}
\usepackage{xcolor}
\usepackage[colorlinks=true, urlcolor=blue!60!black, linkcolor=blue!60!black]{hyperref}
\setlength{\parindent}{0em}
\setlength{\parskip}{0.45em}
\renewcommand{\baselinestretch}{1.0}
\begin{document}
\begin{center}
{\Large \textit{BMW of North America, Inc. v. Gore}, 517 U.S. 559 (1996)\par}
\vspace{0.3em}
{\normalsize Case Brief --- One-Page Edition\par}
\end{center}

\textbf{Memory Jogger.} The repainted-BMW case: the first modern holding that a punitive award was ``grossly excessive'' under due process, with three guideposts (\href{https://www.courtlistener.com/opinion/118026/bmw-of-north-america-inc-v-gore/}{517 U.S. 559}).

\textbf{Facts.} BMW sold cars as new without disclosing repairs costing under 3\% of MSRP. Dr. Gore's repainted sedan lost \$4,000 in value; BMW had sold 983 such cars nationwide, 14 in Alabama.

\textbf{Procedural History.} An Alabama jury awarded \$4 million in punitive damages; the Alabama Supreme Court remitted it to \$2 million; the Supreme Court granted certiorari.

\textbf{Votes.} Reversed, 5--4. Stevens wrote for the Court; Breyer concurred; Scalia and Ginsburg dissented.

\textbf{Holding.} The \$2 million award violated due process. A State may not punish lawful out-of-state conduct, and excessiveness turns on reprehensibility, the ratio of harm to award, and penalties for comparable misconduct.

\textbf{Reasoning.} Reprehensibility was low: the harm was purely economic, safety was never at risk, and BMW could reasonably read state disclosure statutes as safe harbors. The 500:1 ratio was extreme, and Alabama's maximum civil penalty was \$2,000, so BMW lacked fair notice. Breyer stressed Alabama's lack of constraining standards; Scalia called the guideposts ``a road to nowhere''; Ginsburg would have left reform to the States.

\textbf{Examples.} \emph{Same-Side:} \$15 million on a \$10,000 insurance underpayment. \emph{Opposite-Side:} a 4:1 award against a drugmaker hiding fatal side effects. \emph{Fence-Sitter:} a 300:1 award for sabotage threatening far greater potential harm.

\textbf{Critique.} Originalists see revived \textit{Lochner}-style substantive due process; progressives say single-plaintiff ratios underdeter widespread fraud.

\textbf{Key Quotation.} ``Elementary notions of fairness enshrined in our constitutional jurisprudence dictate that a person receive fair notice not only of the conduct that will subject him to punishment, but also of the severity of the penalty that a State may impose.'' --- Justice Stevens, Opinion of the Court, 517 U.S. at 574.

\end{document}
